\documentclass{article}
\usepackage{amsmath,amssymb,amsthm,hyperref}
\newtheorem{proposition}{Proposition}
\title{A scoped regular-module comparison with Libkind--Myers}
\date{16 November 2026}
\begin{document}\maketitle
\paragraph{Audit and attribution.}
This is an application and boundary check, not a new module theorem. Libkind--Myers, \emph{Towards a double operadic theory of systems}, arXiv:2505.18329v2, Definition 3.4 constructs the loose hom bimodule of any pseudo double category, and Theorem 3.15 collapses it to a right module without changing its carrier. Definition 4.1 adds symmetric monoidal structure. Their initial-process doctrine restricts to homs from a chosen point. The present check supplies the project's input to those constructions and tests what does not follow. The paper sketches parts of the general machinery; the specific action below can be verified directly. No Lean verification or novelty certification is asserted.

\section{Pinned input and reconstruction route}
Read the project v0.2 contract for the governing native-preservation requirement. Then read Sections 2 and 4 of \texttt{One\_relative\_construction\_for\_context\_content\_and\_effects.tex}: the relative pseudopromonad gives a bicategory $B$ and pure pseudofunctor $J:C\to B$, and the relative quintet construction supplies a pseudo double category $D=Q_C(B,J)$. Its loose homs retain the native categories $H(a,b)$ and its tight arrows are those of $C$. A square $h\Rightarrow h'$ with sides $u:a\to a'$ and $v:b\to b'$ is a native cell
\[\alpha:J(v)h\Longrightarrow h'J(u).\]
The foundation proves square pasting and companion recovery. Preserve this full $D$, not merely the horizontal bicategory with its pure marking forgotten. Next read the opening scope and A1--A3 of the project external-framing manuscript: an external parameter action is not a binary tensor of effectful processes or an interchange theorem. Finally compare the paper's Definitions 3.4, 3.13 and Theorem 3.15 at \url{https://arxiv.org/html/2505.18329v2}. Those passages determine exactly what data the regular module adds and what it inherits. The companion README gives repository source paths; the external backup retains SHA-256-pinned source copies. This route supports the action check, not an audit of every project theorem or the latest repository history. Reopen the sources if pure variance, native cell orientation or tensor requirements change.

\section{Direct construction}
Let $D_0=C$ and let $D_1$ be the category of loose arrows and squares of $D$. Denote source and target by $s,t:D_1\to D_0$.
\begin{proposition}
The carrier $R=D_1$, with interface projection $p=t$, admits the right action
\[R\times_{D_0}D_1\longrightarrow R,\qquad (h:a\to b,k:b\to c)\longmapsto k\circ h.\]
Its action constraints are the horizontal associator and unitors of $D$. It is the collapse of the regular loose hom bimodule of $D$.
\end{proposition}
\begin{proof}
On a compatible pair of squares $(\alpha,\beta)$, the action is horizontal pasting. In the project's native-cell notation it is
\[J(w)kh\xRightarrow{\beta*h}k'J(v)h\xRightarrow{k'*\alpha}k'h'J(u),\]
with structural reassociations. Identity and vertical-composition preservation are the double-category identity and interchange laws. The output interface is $c$, and a square's output interface map is $w$, so the action is over $D_0$. Naturality of the action associator and unitor is horizontal coherence; their pentagon and triangle are precisely those of $D$. Before collapse the same carrier has projection $(s,t)$ and also the precomposition action. Collapse discards the left action, not the objects or morphisms of $D_1$.
\end{proof}

Thus every native globular cell survives, with non-normal pure units accounted for by canonical unit constraints, and every changing-endpoint square survives. There is no quotient by an effect distribution or coefficient shadow. The target $D$ still contains $C$ and the pure companions $J(u)$; retain their designated choices for pure-compatible comparisons. The collapsed module's interface projection exposes only the output endpoint. Its carrier still has the original source data as part of its displayed objects and squares, but an arbitrary abstract right module does not specify a left endpoint projection or a left action. Keep $(s,t)$ and the regular bimodule when the two-ended interface is required. The original separate context/content presentations and crossing likewise remain marked source data; the claim is not recovery of those presentations from an arbitrary unmarked module.

\section{The fixed-source alternative}
For an object $a$, let $R_a$ have objects $h:a\to b$ and squares whose left tight boundary is $1_a$. Postcomposition defines a right module by the same proof. It is a restriction of the hom bimodule, not the collapsed all-source carrier. It excludes arbitrary-input processes. For a simple witness, the free symmetric monoidal category on one object $A$ has an identity $A\to A$ but no map $I\to A$; fixed-unit initial processes cannot represent that identity by a state at interface $A$. This does not exclude richer encodings. Recovering arbitrary $A\to B$ as initial processes at a changed interface requires an additional representation theorem, e.g. suitable duality; it is not a consequence of the module action. Even when a point is chosen, monoidal closure of its fixed-source carrier requires extra compatible point/tensor structure; the unit-point construction is distinct from an arbitrary fixed point.

\section{Symmetric monoidality is conditional}
If $D$ is supplied with a symmetric monoidal pseudo double structure, then the regular carrier inherits its tensor on loose arrows and squares; $p$ and the action have the same monoidal comparison constraints as target and horizontal composition in $D$. The required compatibility is
\[(kh)\otimes(k'h')\ \cong\ (k\otimes k')(h\otimes h').\]
The monoidal module axioms then follow from the monoidal double-category axioms. This is the monoidal instance of the established hom-and-collapse machinery. The original relative pseudopromonad by itself does not supply that structure. In particular an external frame action $X\times(-)$ and an internal coefficient tensor are not replacements for this displayed effectful interchange.

There is no general automatic upgrade: take the base category terminal and $H(*,*)$ the discrete category on the free noncommutative monoid on $a,b$, with multiplication concatenation. This is a relative pseudopromonad and yields the required ordinary double category and regular module. A monoidal structure on that same one-object horizontal category would make endomorphisms of its sole object (the monoidal unit) commutative by the Eckmann--Hilton argument, contradicting $ab\ne ba$. This is a counterexample in the general relative ambient, not a claim that every O1 doctrine or the particular finite-measure writer class fails to be monoidal.

\section{No automatic joint assembly or feedback}
The action uses exactly existing composition. It neither constructs an admissible coupling nor supplies products of parallel native processes. For example, a discrete hom category with two distinct objects has no product of those objects, while its regular module can still exist. Any application to shared interaction must supply the intended joint object and its projections/comparison maps, and prove their compatibility with the action. Rebranding all effectful processes as interactions does not separate a genuinely independent wiring theory from systems; a further restriction/interpretation is needed for that purpose.

Even a symmetric monoidal regular module need not admit trace. Let $F$ be the free symmetric monoidal category on distinct objects $A,B,X$ and a generator $f:A\otimes X\to B\otimes X$. There is no arrow $A\to B$: a diagram with input $A$ has no $X$ wire, and neither symmetries nor identities create one, so $f$ cannot occur. A total trace would have to send $f$ to such an arrow. View $F$ as a locally discrete bicategory and use the symmetric monoidal structural subgroupoid as pure base (identities, associators, unitors and symmetries); its symmetric monoidal relative quintet double category nevertheless has the regular module above. Thus feedback is additional structure, not generated by the action. This does not obstruct the project's particular partial, retained-state, guarded or finite-measure feedback operators. Transporting those operators requires their own native-cell and action compatibility checks.

\paragraph{Conclusion and next obligation.}
The unmonoidal regular-module connection is an immediate, faithful repackaging at the carrier level using already established constructions on both sides. A full Libkind--Myers module of systems follows conditionally from a verified symmetric monoidal double structure. No new feedback or joint-interaction theorem has been obtained. A substantive next bridge would specify an independent double category of admitted interactions and interpret its action on native processes, proving that selected shared-law assemblies and a selected feedback semantics agree with that action. That stronger statement must not be confused with the regular construction.
\end{document}
